\begin{lemma}[Super-sequence badness gives base badness]
If $f:F\to Q$ is a bad super-sequence on a front on $X$, then there is a
base multi-sequence $h:\mathsf{BaseIncSeq}(X)\to Q$ which is locally constant
and bad, and which satisfies the extension value equation
\[
\forall x\in\mathsf{BaseIncSeq}(X)\;\exists s\in F\;\bigl(
  s\segs\operatorname{range}(x)\mathbin{\wedge}h(x)=f(s)\bigr).
\]
\end{lemma}
\begin{proof}
The structure \noderef{SuperSequence} records both the front and the value
map.  Write $h_F=f.\mathsf{front}$ and $h_f=f.\mathsf{value}$.  Apply
\noderef{SuperSequenceExtension} to $F$, $X$, $h_F$, and $f$.  We obtain a
base multi-sequence $h$, a proof
$h_{\mathrm{lc}}:\mathsf{BaseLocallyConstant}(h)$, and the extension
equation
\[
 \forall y\in\mathsf{BaseIncSeq}(X),\quad
 \exists u\in\operatorname{Finset}(\mathbb N)\;\exists p:
 \mathsf{FrontPrefix}(F,u,y.1),\quad
 h(y)=h_f\!\left(\langle u,p.1\rangle\right).              \tag{1}
\]
The type of $h$ is the base multi-sequence type, and the first component of
the conclusion is the required local-constancy assertion by
\noderef{BaseLocallyConstant}.

We prove the remaining badness assertion.  Fix an arbitrary
$x\in\mathsf{BaseIncSeq}(X)$.  By \noderef{BaseShift}, the shifted object
$x':=\mathsf{BaseShift}_X(x)$ is again an element of
$\mathsf{BaseIncSeq}(X)$, and its first component satisfies
\[
 x'.1=(\mathsf{BaseShift}_X(x)).1=\mathsf{ShiftMap}(x.1).       \tag{2}
\]
Apply (1) first to $x$ and choose witnesses $s$ and
$p_s:\mathsf{FrontPrefix}(F,s,x.1)$, together with
\[
 e_s:\qquad h(x)=h_f\!\left(\langle s,p_s.1\rangle\right).    \tag{3}
\]
Apply (1) again to $x'$ and choose witnesses $t$ and
$p_t:\mathsf{FrontPrefix}(F,t,x'.1)$, together with
\[
 e_t:\qquad h(x')=h_f\!\left(\langle t,p_t.1\rangle\right).  \tag{4}
\]

Unpack the two \noderef{FrontPrefix} witnesses.  Using the definition of
\noderef{ProperPrefixSet}, their first components give
\[
 s\in F,\qquad
 \mathsf{ProperPrefixSet}\bigl(s,\operatorname{range}(x.1)\bigr), \tag{5}
\]
and
\[
 t\in F,\qquad
 \mathsf{ProperPrefixSet}\bigl(t,\operatorname{range}(x'.1)\bigr).
                                                                    \tag{6}
\]
By (2), the second relation in (6) is exactly
\[
 \mathsf{ProperPrefixSet}\bigl(t,\operatorname{range}(\mathsf{ShiftMap}(x.1))\bigr).
                                                                    \tag{7}
\]
Here $x.1$ is an increasing enumeration because $x$ has the subtype type
\mathsf{BaseIncSeq}(X), as defined in \noderef{BaseIncSeq}.

We now build the precise witness required by \noderef{FiniteShift}.  Use
the increasing enumeration $x.1$ as its witness.  Relation (5) is the
first conjunct in the definition of \mathsf{FiniteShift}, and (7) is the
second conjunct, with the shifted enumeration written as
\mathsf{ShiftMap}(x.1).  Therefore
\[
 \mathsf{FiniteShift}(s,t)                                      \tag{8}
\]
holds.  The membership facts $s\in F$ and $t\in F$ are precisely the
first components of $p_s$ and $p_t$.

By the definition of \noderef{BadSuperSequence}, apply the assumed
badness of $f$ to $s$, $t$, these two membership proofs, and the finite
shift witness (8).  We obtain
\[
 \neg r\!\left(h_f\!\left(\langle s,p_s.1\rangle\right),
                 h_f\!\left(\langle t,p_t.1\rangle\right)\right). \tag{9}
\]
Rewrite the two arguments in (9) using the extension equations (3) and
(4).  Since $x'=\mathsf{BaseShift}_X(x)$, this gives
\[
 \neg r\bigl(h(x),h(\mathsf{BaseShift}_X(x))\bigr).             \tag{10}
\]
The element $x\in\mathsf{BaseIncSeq}(X)$ was arbitrary, so (10) holds for
every such $x$.  By the definition of \noderef{BaseBad}, this is exactly
$\mathsf{BaseBad}(r,h)$.  Together with $h_{\mathrm{lc}}$, the extension
equation (1), and the construction of $h$, it proves the stated
conclusion.
\end{proof}
